\label{app:review-example}

This appendix reproduces excerpts from the baseline and graph-augmented
reviews of \texttt{click} pull request \#3493, which narrows the
\texttt{echo} function's \texttt{message} parameter from
\texttt{typing.Any} to \texttt{object} and replaces the type-dispatch
chain with a \texttt{match} statement.  The two reviews were generated by
GPT-4o under matched conditions; the prompts differ only in their
references to the available context, and the graph-augmented arm
additionally receives the structural block.
Sections are condensed for space; the full reviews are archived in the
study repository.

\bigskip

\noindent
\begin{minipage}[t]{0.48\textwidth}
\noindent\textbf{Baseline review} (diff only)\par
\smallskip\hrule\smallskip
{\small
\textbf{Problems found}
\begin{enumerate}[leftmargin=*,itemsep=2pt]
\item The change from \texttt{t.Any} to \texttt{object} may break callers
      that rely on more permissive typing.
\item The \texttt{match} statement may not be compatible with older Python
      versions.
\item Handling of \texttt{None} as a message is not explicitly tested.
\item The interaction of \texttt{nl} with \texttt{None} and empty strings
      is not thoroughly tested.
\end{enumerate}

\textbf{Files referenced}
\begin{itemize}[leftmargin=*,itemsep=1pt]
\item \texttt{src/click/utils.py:236}
\item \texttt{src/click/utils.py:287--295}
\item \texttt{tests/test\_utils.py:98--102}
\end{itemize}
}
\end{minipage}%
\hfill
\begin{minipage}[t]{0.48\textwidth}
\noindent\textbf{Graph-augmented review} (diff + KG context)\par
\smallskip\hrule\smallskip
{\small
\textbf{Problems found}
\begin{enumerate}[leftmargin=*,itemsep=2pt]
\item The type change from \texttt{t.Any} to \texttt{object} might break
      callers that pass types not explicitly handled.
\item The change could violate existing API contracts if callers expect
      \texttt{echo} to accept additional types.
\item \texttt{echo} is widely used across the codebase, including in
      \texttt{src/click/core.py} and \texttt{src/click/termui.py}, which
      could lead to unexpected behaviour.\label{itm:kg-callers}
\item No test covers \texttt{echo} with complex objects that are not
      \texttt{str}, \texttt{bytes}, or \texttt{bytearray}.
\end{enumerate}

\textbf{Files referenced}
\begin{itemize}[leftmargin=*,itemsep=1pt]
\item \texttt{src/click/utils.py:232--311}
\item \texttt{tests/test\_utils.py:98--104}
\item \texttt{src/click/core.py}\;\textsuperscript{$\star$}
\item \texttt{src/click/termui.py}\;\textsuperscript{$\star$}
\item \texttt{src/click/testing.py}\;\textsuperscript{$\star$}
\end{itemize}

\smallskip
{\footnotesize $\star$\;File absent from the diff; surfaced by the
knowledge graph as dependent on the changed module.}
}
\end{minipage}

\bigskip

The baseline review identifies risks that are visible in the diff---the
type narrowing and the new \texttt{match} syntax---but cannot name
\emph{where} those risks materialise beyond the changed file itself.  The
graph-augmented review names three additional files
(\texttt{core.py}, \texttt{termui.py}, \texttt{testing.py}) that depend on
the changed module, two of which call \texttt{echo} at runtime.  The
study's answer key confirms that
\texttt{core.py} and \texttt{termui.py} contain direct call sites of the
changed function; \texttt{testing.py} imports utilities from the same
module.  This is the pattern the primary endpoint (H1) was designed to
detect: the graph supplies evidence that lets the reviewer ground a claim
about affected code in a concrete file reference rather than a generic
warning.
